\documentclass[12pt]{article}
\usepackage[utf8]{inputenc}
\usepackage[T1]{fontenc}
\usepackage{hyperref}
\title{Rjili Language Resources for K-12 Curriculum Development}
\author{Compiled by Claude Sonnet 5 (claude-sonnet-5) for Andr\'as Kornai}
\date{}
\begin{document}
\maketitle
\section{Introduction}
Rjili (also Jijili, Tanjijili, Ujijili; ISO 639-3 uji) is a Plateau language of Niger State, Nigeria, spoken in the eastern part of Paikoro Local Government Area. It is one of several languages known by the ethnic name Koro and should not be confused with Jili (Lijili, mgi). Speaker counts vary by source (8,500 in an older Wikipedia figure, 14,000 per Global Recordings Network, 17,000 per Joshua Project). Search results state that Hausa is the language of wider communication and that Rjili is used in some villages for the first two years of primary school. This is the only evidence of school use that I found, and I found no curriculum materials.
\section{Primary Educational Resources}
No primers, textbooks or curricula were found in the search results.
\section{Linguistic and Reference Materials}
\begin{itemize}
\item Wikipedia, Jijili language: \url{https://en.wikipedia.org/wiki/Jijili_language}
\item Wikipedia, Jili language (a different language, listed to avoid confusion): \url{https://en.wikipedia.org/wiki/Jili_language}
\end{itemize}
\section{Religious and Community Resources}
\begin{itemize}
\item Joshua Project profile of Koro, Jijili in Nigeria: \url{https://joshuaproject.net/people_groups/12798/NI}
\end{itemize}
\section{Digital and Audio Resources}
\begin{itemize}
\item Global Recordings Network page for Rjili (uji): \url{https://globalrecordings.net/en/language/uji}
\end{itemize}
\section{Curriculum Development Context}
Very little is documented for Rjili. Any curriculum would need new community-based work on orthography and materials, with Hausa as the bridge language.
\end{document}